\section{Evidence-Grounded Persona Construction}
\label{sec:pa-personas}

The first stage of the workflow turns the product idea into a set of interviewable personas in four steps: the planner writes an elicitation plan, the corpus builder collects real documents of the domain, the stakeholder miner extracts what stakeholder groups say in them, and the persona generator builds the personas from these observations. The pipeline adapts the persona construction of MAJ-EVAL~\cite{chen2026majeval}, which derives stakeholder groups and their evaluative dimensions from domain literature, from evaluation to requirements elicitation. Figure~\ref{fig:pa-grounding} shows the activities and the objects they exchange; the trace links between these objects are shown in Figure~\ref{fig:pa-trace-model}.

\begin{figure}[H]
    \centering
    \resizebox{\linewidth}{!}{% UML activity diagram with object flows of the persona construction (Section 4.3).
% The trace links between the objects are shown in the class diagram of the trace model.
\begin{tikzpicture}[uml, action/.append style={text width=2.2cm, minimum width=2.4cm}, objnode/.append style={text width=2.5cm, font=\scriptsize}]
    \node[initial] (init) at (-1.7,0) {};
    \node[action] (plan) at (0,0) {Plan the elicitation};
    \node[objnode] (oplan) at (3.2,0) {\underline{:ElicitationPlan}\\\tiny stakeholder types, topics};
    \node[action] (corpus) at (6.4,0) {Build the research corpus};
    \node[objnode] (osrc) at (9.6,0) {\underline{:SourceDocument}\\\tiny [\texttt{target\_sources\_num}]};
    \node[action] (mine) at (9.6,-2.0) {Mine each document};
    \node[objnode] (oobs) at (6.4,-2.0) {\underline{:StakeholderObservation}\\\tiny [*] dimension, evidence};
    \node[action] (gen) at (3.2,-2.0) {Cluster and build the personas};
    \node[objnode] (oper) at (0,-2.0) {\underline{:Persona}\\\tiny [\texttt{target\_personas\_num}]};
    \node[final] (fin) at (-1.7,-2.0) {};
    \node[objnode, text width=2.1cm] (up) at (5.0,1.6) {\underline{:UploadedDocument}\\\tiny [*] from members};
    \node[objnode, text width=2.1cm] (hits) at (7.8,1.6) {\underline{:SearchHit}\\\tiny [*] web, OpenAlex, arXiv};

    \draw[flow] (init) -- (plan);
    \draw[flow] (plan) -- (oplan);
    \draw[flow] (oplan) -- (corpus);
    \draw[flow] (up) -- (corpus);
    \draw[flow] (hits) -- (corpus);
    \draw[flow] (corpus) -- (osrc);
    \draw[flow] (osrc) -- (mine);
    \draw[flow] (mine) -- (oobs);
    \draw[flow] (oobs) -- (gen);
    \draw[flow] (gen) -- (oper);
    \draw[flow] (oper) -- (fin);
\end{tikzpicture}
}
    \caption{UML activity diagram with object flows of the persona construction.}
    \label{fig:pa-grounding}
\end{figure}

\subsection{Elicitation Planning}
\label{sec:pa-planning}

The planner is responsible for deciding what the elicitation must find out. It reads the product idea and the answers the user gave to the clarifying questions asked before the planning, and writes the \emph{elicitation plan}: the objectives of the elicitation, the topics to cover, the stopping conditions, and the stakeholder types that the interviews must represent, such as a student or a member of the parking office staff. These stakeholder types become a closed vocabulary: every persona built later must declare one of them as its stakeholder type, which the system checks.

\subsection{Research Corpus}
\label{sec:pa-corpus}

The corpus builder is responsible for the \emph{research corpus}, the set of documents in which real stakeholders speak. Its inputs are the plan, the documents uploaded by the project members, which come first, and the results of its own web searches: it is the only role with a tool, a search that it calls repeatedly with queries of its own within a bounded number of rounds, until the stakeholder types of the plan are covered by documents of different kinds. From these candidates it keeps the requested number of documents that report what stakeholders do, need, fear or value, preferring studies, surveys, interviews and forum threads to vendor pages, and records for each one an excerpt and the reason it is relevant. The agent can only keep a document that the search or the members actually provided: an output that names any other document is rejected, and a run that finds no usable document stops with an explicit error instead of falling back to invented personas.

\subsection{Stakeholder Observation Mining}
\label{sec:pa-mining}

The stakeholder miner is responsible for extracting, from one document at a time, what the stakeholders in it care about. Each \emph{stakeholder observation} it returns records the group of people concerned, what they say or do, the evaluative dimension that this reveals (for example, the time spent finding a parking spot), the sentence that serves as evidence, and the document that contains it. Only the voice of stakeholders is recorded, never the claims of the authors or the marketing of vendors. The evidence must match the document word for word, which the system checks; observations are therefore the verified ground on which the personas are built, linked to their documents by \emph{evidenced by} links.

\subsection{Persona Generation}
\label{sec:pa-persona-generation}

The persona generator is responsible for turning the observations into the requested number of personas. It groups the observations that describe the same real group, merges those that evaluate the same thing, and builds one persona per distinct evaluative dimension or set of dimensions that one person would carry together. Each persona has the attributes recommended for RE personas~\cite{karolita2023what}, among them goals, constraints, priorities and concerns, and a stakeholder type taken from the vocabulary of the plan; it is linked by \emph{grounded in} links to the observations it is built from, together with their quotes. Personas are given conflicting priorities where real stakeholders would have them, so that the interviews can surface conflicts.

For example, from a campus parking corpus, the observation ``Commuters report circling lots for 10 to 15 minutes before giving up and parking off campus'' of the group \emph{commuter students} leads to the persona ``Maya, 21, third-year biology major commuting 35 minutes'', of stakeholder type \emph{Student}, whose evaluative dimension is the time spent finding a spot.

The run then pauses at the persona checkpoint, where the user can edit the personas before approving them (Section~\ref{sec:pa-checkpoints}).
